\documentclass[
	% -- opções da classe memoir --
	12pt,				% tamanho da fonte
	openright,			% capítulos começam em pág ímpar (insere página vazia caso preciso)
	oneside,			% para impressão em verso e anverso. Oposto a oneside
	a4paper,			% tamanho do papel. 
	% -- opções da classe abntex2 --
	chapter=TITLE,		% títulos de capítulos convertidos em letras maiúsculas
	section=TITLE,		% títulos de seções convertidos em letras maiúsculas
	%subsection=TITLE,	% títulos de subseções convertidos em letras maiúsculas
	%subsubsection=TITLE,% títulos de subsubseções convertidos em letras maiúsculas
	% -- opções do pacote babel --
	english,			% idioma adicional para hifenização
	brazil				% o último idioma é o principal do documento
	]{abntex2}
	
% ---
% Pacotes básicos 
% ---
\usepackage{times}			    % Usa a fonte ADOBE Times			
\usepackage[T1]{fontenc}		% Selecao de codigos de fonte.
\usepackage[utf8]{inputenc}		% Codificacao do documento (conversão automática dos acentos)
\usepackage{lastpage}			% Usado pela Ficha catalográfica
\usepackage{indentfirst}		% Indenta o primeiro parágrafo de cada seção.
\usepackage{color}				% Controle das cores
\usepackage{graphicx}			% Inclusão de gráficos
\usepackage{microtype} 			% para melhorias de justificação
\usepackage{mathtools}          % para uso de funções matemáticas
\usepackage{xfrac}              % para usar frações com \
%\usepackage{multirow}           % Usado para mesclar linhas e colunas de uma tabela
% ---	


% ---
% Pacotes de citações
% ---
%\usepackage[brazilian,hyperpageref]{backref}	 % Paginas com as citações na bibl
\usepackage[alf,bibjustif]{abntex2cite}	% Citações padrão ABNT

% --- 
% CONFIGURAÇÕES DE PACOTES
% --- 

% ---
% Configurações do pacote backref
% Usado sem a opção hyperpageref de backref
%\renewcommand{\backrefpagesname}{Citado na(s) página(s):~}
%% Texto padrão antes do número das páginas
%\renewcommand{\backref}{}
%% Define os textos da citação
%\renewcommand*{\backrefalt}[4]{
%	\ifcase #1 %
%		Nenhuma citação no texto.%
%	\or
%		Citado na página #2.%
%	\else
%		Citado #1 vezes nas páginas #2.%
%	\fi}%
% ---

% Fontes dos títulos e seções (no texto e no sumário, respectivamente)
\renewcommand{\ABNTEXchapterfont}{\bfseries}
\renewcommand{\ABNTEXchapterfontsize}{\normalsize}
%
\renewcommand{\ABNTEXsectionfont}{\normalfont}
\renewcommand{\ABNTEXsectionfontsize}{\normalsize}
%
\renewcommand{\ABNTEXsubsectionfont}{\bfseries}
\renewcommand{\ABNTEXsubsectionfontsize}{\normalsize}
%
\renewcommand{\ABNTEXsubsubsectionfont}{\normalfont}
\renewcommand{\ABNTEXsubsubsectionfontsize}{\normalsize}

%
\renewcommand{\cftchapterfont}{\normalfont\bfseries}
\renewcommand{\cftchapterpagefont}{\normalsize}
%
\renewcommand{\cftsectionfont}{\normalfont\uppercase}
\renewcommand{\cftsectionpagefont}{\normalsize}
%
\renewcommand{\cftsubsectionfont}{\normalsize\bfseries}
\renewcommand{\cftsubsectionpagefont}{\normalsize}
%
\renewcommand{\cftsubsubsectionfont}{\normalfont}
\renewcommand{\cftsubsubsectionpagefont}{\normalsize}
  
%Substitui o texto "Referências" por "REFERÊNCIAS" no sumário
\addto\captionsbrazil{%
  \renewcommand{\bibname}%
    {\uppercase{Referências}}%
}    

% ---
% Informações de dados para CAPA e FOLHA DE ROSTO
% ---
\titulo{Seu título}
\autor{Seu nome}
\local{Cidade}
\data{Ano}
\orientador{Prof. Fulano, Dr. Eng.}
\coorientador{Prof. Beltrano, Dr.}
\instituicao{Universidade}
\tipotrabalho{Dissertação (Mestrado)}
% O preambulo deve conter o tipo do trabalho, o objetivo, 
% o nome da instituição e a área de concentração 
\preambulo{Dissertação apresentada ao Programa de x da Universidade y, como requisito para obtenção do grau de Mestre em z.}
% ---

% ---
% Configurações de aparência do PDF final

% alterando o aspecto da cor azul
\definecolor{blue}{RGB}{41,5,195}

% informações do PDF
\makeatletter
\hypersetup{
     	%pagebackref=true,
		pdftitle={\@title}, 
		pdfauthor={\@author},
    	pdfsubject={\imprimirpreambulo},
	    pdfcreator={LaTeX with abnTeX2},
%	    Coloque suas keyword aqui!!
		pdfkeywords={keyword 1}{keyword 2}{keyword 3}{keyword 4}, 
		colorlinks=true,       		% false: boxed links; true: colored links
    	linkcolor=black,          	% color of internal links
    	citecolor=black,        		% color of links to bibliography
    	filecolor=black,      		% color of file links
		urlcolor=black,
		bookmarksdepth=4
}
\makeatother
% --- 

% --- 
% Espaçamentos entre linhas e parágrafos 
% --- 

% O tamanho do parágrafo é dado por:
\setlength{\parindent}{1.3cm}

% Controle do espaçamento entre um parágrafo e outro:
\setlength{\parskip}{0.2cm}  % tente também \onelineskip

% Cotrole do espaçamento entre o título de um capítulo e o ínicio do seu texto referente
% NBR 14724:2011 - "Os títulos das seções primárias devem começar em página ímpar (anverso), na parte superior da mancha gráfica e ser separados do texto que os sucede por um espaço entre as linhas de 1,5."
\setlength\afterchapskip{18pt}

\setlength\aftersecskip{18pt}

\setlength\aftersubsecskip{18pt}

\setlength\aftersubsubsecskip{18pt}

% ---
% compila o indice
% ---
\makeindex
% ---
	
\begin{document}

% Retira espaço extra obsoleto entre as frases.
\frenchspacing 

% ----------------------------------------------------------
% ELEMENTOS PRÉ-TEXTUAIS
% ----------------------------------------------------------
\pretextual

% ---
% Capa
% ---
\renewcommand{\imprimircapa}{
 \begin{capa}
  \center
  \includegraphics[scale=1]{Brasao_UFBA.png}  \\
  {\bfseries\Large\MakeUppercase\imprimirinstituicao}
   \\
  {\bfseries\large\MakeUppercase{Programa x}}
  \\
  {\bfseries\large\MakeUppercase{Mestrado em x}}
   \vfill
 {\bfseries\large\MakeUppercase\imprimirautor}
 \vfill
 {\bfseries\Large\MakeUppercase\imprimirtitulo}
 \vfill
 {\large\imprimirlocal}
 \par
 {\large\imprimirdata}
 \vspace*{1cm}
 \end{capa}
 }
\imprimircapa
% ---

% ---
% Folha de rosto
% (o * indica que haverá a ficha bibliográfica)
% ---
\makeatletter
\renewcommand{\folhaderostocontent}{
\begin{center}
{\large\MakeUppercase\imprimirautor}
\vspace*{\fill}
\begin{center}
\bfseries\Large\MakeUppercase\imprimirtitulo
\end{center}
\vspace*{\fill}
\abntex@ifnotempty{\imprimirpreambulo}{%
\hspace{.3\textwidth}
\begin{minipage}{.6\textwidth}
\SingleSpacing
\imprimirpreambulo\par
\vspace*{.5cm}
}%
{\imprimirorientadorRotulo~\par\imprimirorientador\par}
\vspace*{.5cm}
{\abntex@ifnotempty\imprimircoorientador\imprimircoorientadorRotulo~\par\imprimircoorientador}
\end{minipage}%
\vspace*{\fill}
{\large\imprimirlocal}
\par
{\large\imprimirdata}
\vspace*{1cm}
\end{center}
}
\makeatother
\imprimirfolhaderosto[Folha de rosto]
% ---

% --- Resumo em Português --- %
\noindent
SOBRENOME, Nome. \imprimirtitulo . Ano. (\imprimirtipotrabalho). Programa de Pós-graduação em x. Universidade y, Cidade, Ano.
\vspace*{1cm}
\begin{resumo}[\large\uppercase{Resumo}]
\vspace*{1cm}
Seu resumo aqui!
 \\
 \\
 \noindent
 \textbf{Palavras-chaves}: keyword 1, keyword 2, keyword 3, keyword 4.
\end{resumo}

% --- Resumo em Inglês --- %
\noindent
SOBRENOME, Nome. \imprimirtitulo . Ano. (Type of work). Programa de Pós-graduação em x. Universidade y, Cidade, Ano.
\vspace*{1cm}
\begin{resumo}[\large\uppercase{Abstract}]
\vspace*{1cm}
Your abstract here!
 \\
 \\
 \noindent
 \textbf{Palavras-chaves}: keyword 1, keyword 2, keyword 3, keyword 4.
\end{resumo}

% --- Sumário --- %
\pdfbookmark[0]{\contentsname}{toc}
\tableofcontents*
\cleardoublepage

% ----------------------------------------------------------
% ELEMENTOS TEXTUAIS
% ----------------------------------------------------------

\textual  %Indica que aqui começam os elementos textuais do documento, isso cria, a partir daqui a numeração de páginas!

%Retira indicação de capítulo e seção no cabeçalho da página deixando apenas a numeração de página.
\pagestyle{simple}

%Sua dissertação aqui!
%Recomendo criar arquivos para cada capítulo e chamá-los aqui assim: 

% ----------------------------------------------------------
% ELEMENTOS PÓS-TEXTUAIS
% ----------------------------------------------------------
\postextual
% ----------------------------------------------------------

% ----------------------------------------------------------
% Referências bibliográficas
% ----------------------------------------------------------
%\bibliography{refs-dissertacao_v1}

\end{document}